\documentclass[conference]{IEEEtran}
\IEEEoverridecommandlockouts

% ── Packages ──────────────────────────────────────────────────────────────────
\usepackage{cite}
\usepackage{amsmath,amssymb,amsfonts}
\usepackage{algorithmic}
\usepackage{algorithm}
\usepackage{graphicx}
\usepackage{textcomp}
\usepackage{xcolor}
\usepackage{booktabs}
\usepackage{multirow}
\usepackage{url}
\usepackage{balance}

\def\BibTeX{{\rm B\kern-.05em{\sc i\kern-.025em b}\kern-.08em
    T\kern-.1667em\lower.7ex\hbox{E}\kern-.125emX}}

% ── Document ──────────────────────────────────────────────────────────────────
\begin{document}

\title{Performance Analysis of Grover's Search Algorithm\\
Under Realistic Noise Models on NISQ Simulators}

\author{
  \IEEEauthorblockN{Saim Haider}
  \IEEEauthorblockA{
    \textit{Dept. of Computer Science}\\
    \textit{FAST-NUCES, Pakistan}\\
    22P-9244\\
    p229244@pwr.nu.edu.pk
  }
  \and
  \IEEEauthorblockN{Muhammad Abdullah}
  \IEEEauthorblockA{
    \textit{Dept. of Computer Science}\\
    \textit{FAST-NUCES, Pakistan}\\
    22P-9371\\
    p229371@pwr.nu.edu.pk
  }
  \and
  \IEEEauthorblockN{Abdullah}
  \IEEEauthorblockA{
    \textit{Dept. of Computer Science}\\
    \textit{FAST-NUCES, Pakistan}\\
    22P-9358\\
    p229358@pwr.nu.edu.pk
  }
}

\maketitle

% Abstract
\begin{abstract}
Grover's search algorithm offers a quadratic speedup over classical unstructured
search, making it one of the most practically relevant quantum algorithms for
near-term deployment. Current quantum hardware operates in the Noisy
Intermediate-Scale Quantum (NISQ) regime, where decoherence and gate errors
can significantly degrade algorithm performance. This paper examines the behavior
of Grover's algorithm under three realistic noise models, namely depolarizing
noise, bit-flip errors, and phase-flip errors, through a controlled simulation
study. Using Python and the Qiskit Aer simulator, Grover's search circuit is
implemented for search spaces of $n = 2, 3, 4, 5$ qubits, with noise injected
at error rates swept across $p \in [0.001,\, 0.1]$ and success probability
recorded at each configuration. The results indicate that Grover's algorithm
retains quantum advantage ($P_{\text{success}} > 1/N$) across the full tested
error range for all configurations, with the critical threshold $p^{*}$ remaining
above $0.100$ in every case. Phase-flip and bit-flip noise produce nearly
identical degradation patterns, while depolarizing noise shows comparatively
less impact at small qubit counts. Crossover analysis and comparison against
the theoretical predictions of Salas (2008) and Dowarah et al.\ (2025) are
included and discussed.
\end{abstract}

\begin{IEEEkeywords}
Grover's algorithm, NISQ computing, quantum noise models, quantum advantage,
Qiskit Aer, depolarizing noise, bit-flip noise, phase-flip noise
\end{IEEEkeywords}

% ── I. Introduction ───────────────────────────────────────────────────────────
\section{Introduction}

Quantum computing leverages superposition, entanglement, and quantum interference
to solve certain problems significantly faster than classical computers. Among the
most practically relevant quantum algorithms is Grover's search algorithm~\cite{grover1996},
which achieves $\mathcal{O}(\sqrt{N})$ query complexity for searching an unstructured
database of $N$ entries, compared to the classical $\mathcal{O}(N)$ average case.
This quadratic speedup has wide applications in optimization, cryptanalysis, and
database search.

Despite its theoretical promise, quantum algorithms face a fundamental practical
challenge: noise. Physical qubits are susceptible to environmental decoherence,
gate imperfections, and measurement errors. These noise sources corrupt the quantum
interference patterns that Grover's algorithm relies upon to amplify the probability
of measuring the correct answer. Present-day quantum processors operate in the NISQ
regime, characterized by gate error rates on the order of $10^{-3}$ to $10^{-2}$
and coherence times in the microsecond to millisecond range. Under these conditions,
even moderately deep circuits such as Grover's can suffer severe degradation in
success probability.

This work presents a systematic empirical study of Grover's algorithm under
realistic noise conditions using the Qiskit Aer simulator. By sweeping noise error
rates across a controlled range and measuring success probability for each
configuration, this paper identifies the precise noise threshold beyond which
Grover's algorithm loses its quantum advantage over classical search. The full
implementation is provided as a reproducible open-source Qiskit notebook.

% ── II. Literature Review ─────────────────────────────────────────────────────
\section{Literature Review}

\subsection{Ijaz and Faryad (2023) -- Base Paper}

Ijaz and Faryad~\cite{ijaz2023} conducted a comprehensive noise analysis of
Grover's algorithm using Qiskit Aer for $n = 3, 4, 5$ qubits with 10{,}000 shots
per configuration. Four noise channels were studied: depolarizing, bit-flip,
phase-flip, and bit-phase flip. The key finding was that success probability decays
exponentially with error probability $p$ but only linearly with qubit count $n$.
However, the study does not identify the critical threshold $p^{*}$ where quantum
advantage over classical search is lost.

\subsection{Salas (2008)}

Salas~\cite{salas2008} performed the earliest systematic numerical study of noise
in Grover's algorithm using Fortran-based Monte Carlo simulation for $n = 2$ to
$7$ qubits under a depolarizing channel model. The study derived an exponential
damping law and the allowed-error scaling relation
$\varepsilon_{\text{th}}(N) \sim N^{-1.1}$, identifying an absolute threshold of
$\varepsilon > 0.043$ above which the algorithm fails for any qubit count. The
limitation is that only depolarizing noise is studied and the framework predates
modern Qiskit tooling.

\subsection{Dowarah et al.\ (2025)}

Dowarah et al.~\cite{dowarah2025} applied Random Matrix Theory (RMT) to model the
noisy Grover operator as a Floquet unitary, deriving two analytical noise
thresholds: the computing transition at $\delta_{c,\text{comp}} \sim
\mathcal{O}(2^{-L/2})$ and the ergodicity transition at $\delta_{c,\text{gap}}
\sim \mathcal{O}(L^{-1})$. The work is entirely theoretical with no Qiskit
simulation provided.

\subsection{Srivastava et al.\ (2025)}

Srivastava et al.~\cite{srivastava2025} proposed using quantum switches exploiting
indefinite causal order to mitigate depolarizing noise in Grover's algorithm. The
study is entirely analytical and does not identify the threshold $p^{*}$ for
standard Grover circuits without mitigation.

\subsection{Research Gap}

No existing study combines (1) a modern Qiskit simulation framework, (2) multiple
isolated noise types, (3) systematic error rate sweeps, and (4) explicit
identification of the quantum advantage loss threshold $p^{*}$. This work fills
that gap.

% ── III. Methodology ──────────────────────────────────────────────────────────
\section{Methodology}

This section describes the complete implementation pipeline for simulating Grover's
search algorithm under realistic noise conditions. The methodology encompasses
circuit construction, noise model injection, simulation execution, and threshold
identification.

\begin{figure}[htbp]
  \centering
  \includegraphics[width=0.82\columnwidth]{fig_methodology_workflow.png}
  \caption{Proposed Methodology Workflow Block Diagram for the Grover Noise
  Analysis Pipeline, showing the seven-stage flow from input specification through
  circuit construction, noise injection, simulation, probability measurement,
  threshold identification, to final output.}
  \label{fig:workflow}
\end{figure}

\subsection{Circuit Construction}

Grover's search circuit is implemented from first principles in Qiskit for
$n = 2, 3, 4, 5$ qubits. Each circuit consists of three components: (1) a
Hadamard initialization layer placing all qubits into uniform superposition,
(2) an oracle operator $\mathcal{O}_\tau = 2|\tau\rangle\langle\tau| - I$ marking
the target state $|\tau\rangle = |11\ldots1\rangle$, implemented via an
H-MCX-H decomposition, and (3) a Grover diffusion operator $D = 2|s\rangle\langle s|
- I$ performing inversion about the mean. The optimal number of iterations is
$k = \lfloor (\pi/4)\sqrt{N} \rfloor$, where $N = 2^n$.

Table~\ref{tab:circuit} presents the circuit parameters for each qubit
configuration after transpilation to the AerSimulator basis gate set. Circuit depth
grows linearly with $n$, from 6 for $n=2$ to 18 for $n=5$, while gate count grows
from 9 to 57. These increasing depths accumulate more noise per simulation, making
larger qubit counts progressively more sensitive to gate error rates.

\begin{table}[htbp]
  \centering
  \caption{Circuit Parameters vs.\ Qubit Count}
  \label{tab:circuit}
  \begin{tabular}{ccccc c}
    \toprule
    $n$ & $N=2^n$ & $k$ & Depth & Gates & Baseline $1/N$ \\
    \midrule
    2 & 4  & 1 & 6  & 9  & 0.2500 \\
    3 & 8  & 2 & 10 & 21 & 0.1250 \\
    4 & 16 & 3 & 14 & 37 & 0.0625 \\
    5 & 32 & 4 & 18 & 57 & 0.0313 \\
    \bottomrule
  \end{tabular}
\end{table}

\subsection{Noise Model Injection}

Three noise models are injected independently using the Qiskit Aer \texttt{NoiseModel}
API. Each error channel is applied per gate across all single-qubit gates
(\texttt{H, X, Z, U}) and two-qubit gates (\texttt{CX, CCX}).

\textbf{Depolarizing Noise.} A completely mixed state is applied with probability
$p$, implemented via \texttt{depolarizing\_error(p, 1)} for single-qubit gates and
\texttt{depolarizing\_error(p, 2)} for two-qubit gates. This channel simultaneously
corrupts both phase and amplitude information.

\textbf{Bit-Flip Noise.} A Pauli $X$ gate is applied with probability $p$,
implemented via \texttt{pauli\_error([(`X', p), (`I', 1-p)])}. Two-qubit errors are
constructed as the tensor product of two independent single-qubit errors.

\textbf{Phase-Flip Noise.} A Pauli $Z$ gate is applied with probability $p$,
implemented via \texttt{pauli\_error([(`Z', p), (`I', 1-p)])}. Phase-flip noise
directly disrupts the quantum interference mechanism underlying Grover's amplitude
amplification.

Error rates are swept across $p \in [0.001, 0.1]$ in 20 uniform steps
($\Delta p \approx 0.005$), yielding 1{,}200 total simulation configurations
($3\ \text{noise types} \times 4\ \text{qubit counts} \times 20\ \text{error rates}
\times 5\ \text{runs}$).

\subsection{Simulation Execution}

All simulations are executed on the Qiskit \texttt{AerSimulator} with a fixed
random seed (\texttt{seed = 42}) to ensure full reproducibility. Each configuration
is executed with \texttt{SHOTS = 1024} measurements averaged over \texttt{RUNS = 5}
independent trials. Success probability is computed as:
\begin{equation}
  P_{\text{success}} = \frac{\text{count}(|\tau\rangle)}{1024}
  \label{eq:psuccess}
\end{equation}
averaged over five runs, with error bars representing $\pm 1$ standard deviation.
Circuits are transpiled to the AerSimulator basis gate set prior to noise injection.

\subsection{Threshold Identification}

The critical threshold $p^{*}$ is defined as the error rate at which Grover's
measured success probability crosses the classical random search baseline
$1/N = 1/2^n$. The crossover is identified via linear interpolation:
\begin{equation}
  p^{*} = p_i + \frac{\tfrac{1}{N} - P_{\text{mean}}(p_i)}
  {P_{\text{mean}}(p_{i+1}) - P_{\text{mean}}(p_i)}
  \cdot (p_{i+1} - p_i)
  \label{eq:threshold}
\end{equation}
where $P_{\text{mean}}(p_i) \geq 1/N > P_{\text{mean}}(p_{i+1})$. If no crossover
occurs within the tested range, the threshold is recorded as $p^{*} > 0.100$.

\begin{algorithm}
\caption{Grover Noise Analysis Simulation Pipeline}
\label{alg:pipeline}
\begin{algorithmic}[1]
\REQUIRE $n \in \{2,3,4,5\}$, noise\_type, $p \in [0.001, 0.1]$
\ENSURE Threshold table, success probability curves
\STATE Set SHOTS $\leftarrow$ 1024, RUNS $\leftarrow$ 5, seed $\leftarrow$ 42
\FOR{each $n \in \{2,3,4,5\}$}
  \STATE $k \leftarrow \lfloor (\pi/4)\sqrt{2^n} \rfloor$
  \STATE Build Grover circuit $QC(n,\, k,\, \text{target}=|11\ldots1\rangle)$
  \STATE Transpile $QC$ to AerSimulator basis gate set
  \FOR{each noise\_type}
    \FOR{each $p$ in ERROR\_RATES}
      \STATE $NM \leftarrow$ NoiseModel(noise\_type, $p$)
      \STATE run\_probs $\leftarrow [\,]$
      \FOR{$r = 1$ \TO RUNS}
        \STATE counts $\leftarrow$ AerSimulator.run($QC$, $NM$, SHOTS)
        \STATE run\_probs.append(counts$[|\tau\rangle]$ / SHOTS)
      \ENDFOR
      \STATE stats$[$noise\_type$][n][p] \leftarrow \{\mu,\sigma\}$(run\_probs)
    \ENDFOR
    \STATE $p^{*} \leftarrow$ LinearInterpolate(stats, baseline $= 1/2^n$)
  \ENDFOR
\ENDFOR
\STATE Generate Figs.~1--4 and Tables~I--III
\RETURN thresholds, stats
\end{algorithmic}
\end{algorithm}

% ── IV. Results ───────────────────────────────────────────────────────────────
\section{Results}

This section presents the complete experimental results of simulating Grover's
search algorithm under three noise models across $n = 2, 3, 4, 5$ qubits. All
results are derived from 1{,}200 simulation configurations with 1{,}024 shots
averaged over five independent runs.

\subsection{Noiseless Baseline Validation}

Prior to noise injection, the Grover circuit was validated in an ideal noiseless
environment. For all qubit configurations, the measured success probability exceeded
0.97, confirming correct oracle and diffusion operator implementation. Measured
values reached approximately 0.997 for $n=2$ and 0.973 for $n=5$, with small
deviations attributable to finite shot statistics.

\subsection{Success Probability vs.\ Error Rate}

Fig.~\ref{fig:success} presents the success probability curves for all three noise
models across $n = 2$ to $5$ qubits as error rate $p$ increases from 0.001 to 0.1.
Success probability decays monotonically with increasing $p$ for all configurations,
consistent with the exponential degradation law reported by Ijaz and Faryad~\cite{ijaz2023}
and the analytical results of Salas~\cite{salas2008}. Larger qubit counts experience
steeper decay due to increased circuit depth, as more gate operations accumulate noise
per run. For $n=5$, the circuit applies four Grover iterations across 57 gates at
depth 18, compared to one iteration, nine gates, and depth 6 for $n=2$.

At $p = 0.001$, all configurations achieve $P_{\text{success}} > 0.93$. At $p = 0.1$,
depolarizing noise leaves $n=2$ at $P = 0.655$ against a classical baseline of 0.250,
and $n=4$ at $P = 0.255$ against a baseline of 0.063. For $n=5$ under bit-flip and
phase-flip noise, the success probability at $p = 0.1$ approaches approximately 0.09
to 0.12, still above the classical baseline of 0.031.

\begin{figure*}[htbp]
  \centering
  \includegraphics[width=\textwidth]{fig1_success_probability.png}
  \caption{Grover success probability under depolarizing (left), bit-flip (center),
  and phase-flip (right) noise models for $n = 2, 3, 4, 5$ qubits. Dashed horizontal
  lines indicate classical random search baselines $1/N$. Error bars represent
  $\pm 1$ standard deviation over five independent runs.}
  \label{fig:success}
\end{figure*}

\subsection{Noise Model Comparison}

Comparing across noise channels: depolarizing noise causes the least degradation at
small qubit counts ($n = 2, 3$) but produces comparable degradation at larger $n$.
Bit-flip and phase-flip noise produce nearly identical degradation curves across all
configurations, contrary to the initial hypothesis that phase-flip would be more
severe due to Grover's sensitivity to phase interference. This parity arises because
both Pauli $X$ and Pauli $Z$ errors are single-channel errors with equivalent
probability amplitudes. The critical distinction emerges only in circuit architectures
with strongly asymmetric phase sensitivity.

\subsection{Quantum Advantage Loss Thresholds}

Table~\ref{tab:thresholds} presents the quantum advantage loss threshold $p^{*}$ for
each noise type and qubit count. A value of $> 0.100$ indicates that Grover's algorithm
maintains $P_{\text{success}} > 1/N$ throughout the entire tested error range.

\begin{table}[htbp]
  \centering
  \caption{Quantum Advantage Loss Threshold $p^{*}$ Values}
  \label{tab:thresholds}
  \begin{tabular}{lcccc}
    \toprule
    Noise Type     & $n=2$   & $n=3$   & $n=4$   & $n=5$   \\
    \midrule
    Depolarizing   & $>0.100$ & $>0.100$ & $>0.100$ & $>0.100$ \\
    Bit-Flip       & $>0.100$ & $>0.100$ & $>0.100$ & $>0.100$ \\
    Phase-Flip     & $>0.100$ & $>0.100$ & $>0.100$ & $>0.100$ \\
    \bottomrule
  \end{tabular}
\end{table}

The finding that $p^{*} > 0.100$ for all 12 configurations is a significant positive
result. It indicates that Grover's algorithm is remarkably robust within the error
rate range characteristic of current NISQ hardware, where typical two-qubit gate
error rates lie between $10^{-3}$ and $10^{-2}$.

\subsection{Threshold Scaling Analysis}

Fig.~\ref{fig:scaling} presents the threshold scaling plot comparing empirical results
against theoretical predictions. The Salas (2008) law predicts thresholds decreasing
from approximately 0.058 at $n=2$ to 0.006 at $n=5$. The Dowarah et al.\ (2025)
computing transition falls even more steeply. The empirical threshold (flat boundary
at $p^{*} > 0.100$) indicates that the crossover did not occur within the tested
range for any qubit count or noise type.

\begin{figure}[htbp]
  \centering
  \includegraphics[width=\columnwidth]{fig2_threshold_scaling.png}
  \caption{Quantum advantage loss threshold $p^{*}$ vs.\ qubit count $n$.
  Empirical results (flat boundary at $p^{*} > 0.100$) compared with Salas (2008)
  theoretical scaling law $\varepsilon_{\text{th}}(N) \sim N^{-1.1}$ and
  Dowarah et al.\ (2025) computing transition $\delta_{c,\text{comp}} \sim 2^{-L/2}$.}
  \label{fig:scaling}
\end{figure}

\subsection{Crossover Analysis}

Fig.~\ref{fig:crossover} presents grouped bar charts comparing Grover's success
probability against the classical baseline at two representative error rates:
$p \approx 0.011$ (low noise, near current best NISQ hardware) and $p \approx 0.048$
(moderate noise, representative of near-term hardware). At $p = 0.011$, all Grover
configurations substantially exceed the classical baseline, with $P_{\text{success}}$
ranging from 0.75 to 0.95 compared to classical baselines between 0.031 and 0.250.
At $p = 0.048$, quantum advantage is maintained but reduced: $n=5$ configurations
reach $P \approx 0.27$--$0.30$, still well above the $n=5$ baseline of 0.031.

\begin{figure}[htbp]
  \centering
  \includegraphics[width=\columnwidth]{fig3_crossover_comparison.png}
  \caption{Grover success probability vs.\ classical random search baseline at fixed
  error rates $p \approx 0.011$ (left) and $p \approx 0.048$ (right) across all
  noise types and qubit counts. Grey bars indicate the classical baseline $1/N$.}
  \label{fig:crossover}
\end{figure}

% ── V. Evaluation ─────────────────────────────────────────────────────────────
\section{Evaluation of Proposed Work}

\subsection{Comparison with Baseline Methods}

Fig.~\ref{fig:eval} presents the threshold comparison between this work and two
theoretical baselines. The empirical $p^{*}$ values consistently exceed both
Salas (2008) and Dowarah et al.\ (2025) predictions across all qubit counts. This
implies that the standard Grover implementation with optimal $k$ is more
noise-resilient in practice than Monte Carlo estimates or coherent-noise theoretical
models suggest. The discrepancy is most pronounced at $n=5$, where Salas (2008)
predicts a threshold of approximately 0.006 and Dowarah et al.\ predict approximately
0.003, while empirical measurements show $p^{*} > 0.100$ for all three noise channels.

\begin{figure}[htbp]
  \centering
  \includegraphics[width=\columnwidth]{fig4_evaluation_comparison.png}
  \caption{Threshold comparison: this work (empirical, Qiskit 1.x) vs.\
  Salas (2008) analytical scaling law and Dowarah et al.\ (2025) computing
  transition analytical prediction.}
  \label{fig:eval}
\end{figure}

\subsection{Comparative Analysis}

Table~\ref{tab:comparison} summarizes the comparative analysis between this work
and related methods across five dimensions: framework used, noise types studied,
whether $p^{*}$ is identified, and key limitations.

\begin{table}[htbp]
  \centering
  \caption{Comparative Analysis: This Work vs.\ Related Methods}
  \label{tab:comparison}
  \resizebox{\columnwidth}{!}{%
  \begin{tabular}{p{2cm}p{1.8cm}p{1.5cm}p{1cm}p{2.2cm}}
    \toprule
    Method & Framework & Noise Types & $p^{*}$ & Limitation \\
    \midrule
    \textbf{This Work} & Qiskit 1.x Aer & Dep, BF, PF & $>$0.100 all & Small $n$ only \\
    Ijaz \& Faryad (2023) & Qiskit Aer & Dep,BF,PF,BPF & No & No threshold \\
    Salas (2008) & Fortran MC & Dep only & 0.043 & Single channel \\
    Dowarah et al.\ (2025) & Analytical RMT & Coherent & Analytical & No simulation \\
    Srivastava et al.\ (2025) & Analytical & Dep only & No & No Qiskit \\
    \bottomrule
  \end{tabular}}
\end{table}

\subsection{Discussion of Key Findings}

The finding that $p^{*} > 0.100$ for all configurations has direct implications for
NISQ hardware development. Current state-of-the-art superconducting qubit platforms
exhibit two-qubit gate error rates of approximately 0.001 to 0.01, well below the
measured threshold boundary. This indicates that for small search spaces ($n \leq 5$),
Grover's algorithm can be executed with meaningful quantum advantage on current
hardware, provided circuit compilation is optimized to minimize unnecessary gate
count and depth.

The divergence from Salas (2008) predictions at small $n$ may arise from three
sources. First, Salas used Monte Carlo simulation with per-iteration error models
rather than per-gate injection. Second, the Fortran simulation predates modern gate
decomposition optimization. Third, the $n \leq 5$ regime may lie below the
asymptotic region where the $N^{-1.1}$ scaling applies. Future work extending to
$n = 6$--$8$ and implementing error mitigation via the quantum switch framework of
Srivastava et al.\ (2025) is a natural continuation.

% ── VI. Conclusion ────────────────────────────────────────────────────────────
\section{Conclusion}

This paper presented a comprehensive empirical study of Grover's search algorithm
under three realistic noise models using the Qiskit Aer simulator. Through 1{,}200
simulation configurations spanning $n = 2$ to $5$ qubits and error rates
$p \in [0.001, 0.1]$, four primary conclusions are established.

First, Grover's algorithm maintains quantum advantage throughout the entire
NISQ-relevant error range for all tested configurations, with $p^{*} > 0.100$
across all noise types and qubit counts. Second, success probability decays
monotonically and quasi-exponentially with increasing error rate, consistent with
existing literature. Third, bit-flip and phase-flip noise produce comparable
degradation, while depolarizing noise permits greater resilience at small $n$.
Fourth, empirical thresholds substantially exceed Salas (2008) and Dowarah et al.\
(2025) theoretical predictions within the small-$n$ regime, suggesting that standard
Grover circuits are more noise-resilient in practice than prior models indicate.

These results provide concrete, reproducible benchmarks for NISQ hardware calibration
and establish an empirical foundation for future research in quantum algorithm
robustness and error mitigation.

% ── References ────────────────────────────────────────────────────────────────
\balance
\begin{thebibliography}{00}

\bibitem{grover1996}
L.~K. Grover, ``A fast quantum mechanical algorithm for database search,''
in \textit{Proc. 28th Annual ACM Symposium on Theory of Computing (STOC)},
1996, pp. 212--219.

\bibitem{ijaz2023}
M.~A. Ijaz and M.~Faryad, ``Noise analysis of Grover and phase estimation
algorithms implemented as quantum singular value transformations for a small
number of noisy qubits,'' \textit{Scientific Reports}, vol.~13, p.~20144, 2023.

\bibitem{salas2008}
P.~J. Salas, ``Noise effect on Grover algorithm,''
\textit{The European Physical Journal D}, vol.~46, pp.~365--381, 2008.

\bibitem{dowarah2025}
S.~Dowarah, C.~Zhang, V.~Khemani, and M.~Kolodrubetz,
``Phases and phase transition in Grover's algorithm with systematic noise,''
\textit{arXiv preprint arXiv:2406.10344v2}, 2025.

\bibitem{srivastava2025}
S.~Srivastava, A.~K. Pati, I.~Chakrabarty, and S.~Bhattacharya,
``Using quantum switches to mitigate noise in Grover's search algorithm,''
\textit{arXiv preprint arXiv:2401.05866v2}, 2025.

\end{thebibliography}

\end{document}\begin{figure}
    \centering
    \includegraphics[width=0.5\linewidth]{fig_methodology_workflow.png}
    \caption{Enter Caption}
\begin{figure}
        \centering
        \includegraphics[width=0.5\linewidth]{fig1_success_probability.png}
        \caption{Enter Caption}
\begin{figure}
            \centering
            \includegraphics[width=0.5\linewidth]{fig2_threshold_scaling.png}
            \caption{Enter Caption}
\begin{figure}
                \centering
                \includegraphics[width=0.5\linewidth]{fig3_crossover_comparison.png}
                \caption{Enter Caption}
\begin{figure}
                    \centering
                    \includegraphics[width=0.5\linewidth]{fig4_evaluation_comparison.png}
                    \caption{Enter Caption}
                    \label{fig:placeholder}
                \end{figure}
                                \label{fig:placeholder}
            \end{figure}
                        \label{fig:placeholder}
        \end{figure}
                \label{fig:placeholder}
    \end{figure}
        \label{fig:placeholder}
\end{figure}
